\chapter{Energi Inti: Fisi, Fusi, $E=mc^2$}\label{ch:g12:nuclear-energy}

Sebutir peleti bahan bakar sebesar penghapus pensil menutupi kebutuhan listrik
sebuah rumah tangga selama setahun; jasa yang sama dari batu bara menuntut tiga ton.
Di atas sana, Matahari membayar siangnya dengan melenyapkan dirinya --- empat juta ton
dirinya tiap sekon. Bab ini membuka buku besar di balik kedua tawaran itu: massa itu
sendiri adalah \omterm{def:g10:energy-conservation:energy}{energi}, dan inti (\cref{ch:g11:nucleus-radioactivity}) adalah tempat ia diuangkan.

\section{Massa adalah energi}

\begin{theorem}[Kesetaraan massa--energi]\label{thm:g12:nuclear-energy:emc2}\index{kesetaraan massa--energi}
Sebuah benda bermassa $m$, sekadar karena ada, memegang \omterm{def:g10:energy-conservation:energy}{energi}
\[
E = m c^2, \qquad c = \qty{3.00e8}{m/s}.
\]
Setiap kali sebuah sistem kehilangan massa $\Delta m$, ia melepaskan \omterm{def:g10:energy-conservation:energy}{energi}
$\Delta m\, c^2$ --- apa pun mekanismenya.
\end{theorem}
\admitted

\begin{remark}[Einstein, 1905]\label{rem:g12:nuclear-energy:einstein}
Inilah hubungan Einstein, yang diterima saja di sini; dari mana asalnya adalah urusan
bab kerelatifan khusus yang menutup tahun ini
(\cref{ch:g12:special-relativity}). Perhatikanlah kursnya: $c^2 = \qty{9.00e16}{J/kg}$
--- galat pembulatan massa adalah kekayaan \omterm{def:g10:energy-conservation:energy}{energi}.
\end{remark}

\begin{definition}[Elektronvolt dan satuan massa atom]\label{def:g12:nuclear-energy:units}
Dua \omterm{def:g10:orders-of-magnitude:unit}{satuan} yang dijahit untuk intinya.
\emph{Elektronvolt}\index{elektronvolt}, yaitu \omterm{def:g10:energy-conservation:energy}{energi} yang diperoleh sebuah elektron
saat melintasi \qty{1}{V} (\cref{ch:g11:circuits-and-power}):
$\qty{1}{eV} = \qty{1.60e-19}{J}$, $\qty{1}{MeV} = \qty{1.60e-13}{J}$ ---
ikatan kimia memperdagangkan beberapa \unit{eV}, reaksi inti megaelektronvolt.
\emph{Satuan massa atom terpadu}\index{satuan massa atom terpadu}, yaitu
seperdua belas massa sebutir \omterm{def:g11:fundamental-interactions:ladder}{atom} karbon-12: $\qty{1}{u} = \qty{1.66054e-27}{kg}$,
yang setara \omterm{def:g10:energy-conservation:energy}{energi} \qty{931.5}{MeV} --- $\qty{1}{u} = \qty{931.5}{MeV}/c^2$.
\end{definition}

\begin{example}[Kekayaan yang tertidur]\label{ex:g12:nuclear-energy:fortune}
Satu \omterm{def:g10:orders-of-magnitude:unit}{kilogram} apa pun, yang diubah seluruhnya, bernilai $\qty{9.00e16}{J}$ --- tiga
tahun keluaran sebuah pembangkit besar (\qty{1}{GW}). Tidak ada yang memperhatikannya sebelum
1905 karena fisika biasa hampir tidak menyentuh massa: membakar sekilogram batu bara
(\qty{3.0e7}{J}) meringankannya sepertiga mikrogram. Inti, akan kita lihat,
menggeser satu bagian dalam seribu: masih kecil, tetapi sejuta kali kimia.
\end{example}

\section{Defek massa}

\begin{definition}[Defek massa dan energi ikat]\label{def:g12:nuclear-energy:defect}
Timbanglah sebuah inti ${}^{A}_{Z}\mathrm{X}$ (bermassa $m$), lalu timbanglah $Z$ protonnya
dan $A - Z$ neutronnya secara terpisah: bagian-bagiannya \emph{lebih berat} daripada keseluruhannya.
Selisihnya
\[
\Delta m = Z\,m_p + (A - Z)\,m_n - m > 0,
\qquad m_p = \qty{1.00728}{u},\; m_n = \qty{1.00866}{u},
\]
adalah \emph{defek massa}\index{defek massa}; \emph{energi
ikat}\index{energi ikat} $E_b = \Delta m\,c^2$ adalah \omterm{def:g10:energy-conservation:energy}{energi} yang dibutuhkan untuk
mencerai-beraikan intinya menjadi \omterm{def:g11:fundamental-interactions:ladder}{nukleon} bebas --- sama saja, \omterm{def:g10:energy-conservation:energy}{energi} yang dilepaskan
ketika ia dirakit. Inti yang terikat duduk \emph{di bawah} bagian-bagiannya.
\end{definition}

\begin{omfigure}
\begin{tikzpicture}[font=\small]
  \draw[thick] (0,2.1) -- (2.7,2.1) node[right] {$Z$ proton $+$ $(A-Z)$ neutron, bebas};
  \foreach \x in {0.3,0.9,1.5,2.1} \fill[omThm!80] (\x,2.4) circle (2.4pt);
  \foreach \x in {0.6,1.2,1.8,2.4} \fill[omExo] (\x,2.4) circle (2.4pt);
  \draw[thick, omDef] (0,0) -- (2.7,0) node[right] {inti, bermassa $m$};
  \draw[-Stealth, thick] (1.35,2.1) -- (1.35,0) node[midway, right] {$E_b = \Delta m\,c^2$};
\end{tikzpicture}

{\small Sebuah tangga \omterm{def:g10:energy-conservation:energy}{energi}: inti yang terikat terletak $E_b$ di bawah \omterm{def:g11:fundamental-interactions:ladder}{nukleonnya}
yang terpisah --- merakitnya melepaskan $E_b$, membongkarnya menuntut ongkos itu.}
\end{omfigure}

\begin{example}[Helium-4]\label{ex:g12:nuclear-energy:helium}
Inti ${}^{4}_{2}\mathrm{He}$ memiliki $m = \qty{4.00151}{u}$; bagian-bagiannya:
$2 \times 1.00728 + 2 \times 1.00866 = \qty{4.03188}{u}$. Jadi
$\Delta m = \qty{0.03037}{u}$ dan $E_b = 0.03037 \times 931.5 \approx
\qty{28.3}{MeV}$: intinya $0.75\%$ lebih ringan daripada bagian-bagiannya --- yaitu
``satu bagian dalam seribu'' pada \cref{ex:g12:nuclear-energy:fortune}.
\end{example}

\begin{method}[Neraca energi sebuah reaksi inti]\label{met:g12:nuclear-energy:balance}
\begin{enumerate}
  \item Setimbangkanlah persamaannya: jumlah $A$ dan $Z$ kekal (\cref{ch:g11:nucleus-radioactivity}).
  \item Jumlahkanlah massanya sebelum, lalu sesudah, dalam \unit{u} --- pertahankanlah kelima desimalnya.
  \item $\Delta m = m_{\text{sebelum}} - m_{\text{sesudah}}$, lalu
        $E = \Delta m\,(\text{dalam u}) \times \qty{931.5}{MeV}$, yang terbawa pergi sebagai
        \omterm{def:g11:mechanical-energy:kinetic}{energi kinetik} hasilnya; nilai negatif berarti reaksinya harus dibayar.
\end{enumerate}
\end{method}

\section{Kurva yang menjalankan alam semesta}

\begin{definition}[Energi ikat tiap nukleon]\label{def:g12:nuclear-energy:pernucleon}
Kepaduannya dibandingkan secara adil lewat \emph{energi ikat tiap
nukleon}\index{energi ikat tiap nukleon}, $E_b/A$, yaitu ongkos rata-rata
mencabut satu \omterm{def:g11:fundamental-interactions:ladder}{nukleon}: $28.3/4 \approx \qty{7.1}{MeV}$ untuk helium-4.
\end{definition}

\begin{proposition}[Puncak besi, dan dua jalan turun]\label{prop:g12:nuclear-energy:ironpeak}\index{puncak besi}
Ketika dilukiskan terhadap $A$, \omterm{def:g12:nuclear-energy:pernucleon}{energi ikat tiap nukleon} menanjak curam melewati
inti yang ringan, memuncak di dekat besi, ${}^{56}_{26}\mathrm{Fe}$, pada sekitar
\qty{8.8}{MeV} tiap \omterm{def:g11:fundamental-interactions:ladder}{nukleon}, lalu meluncur landai sampai sekitar \qty{7.6}{MeV} pada
uranium: besi adalah inti yang paling erat terikat di alam. Sebuah reaksi
melepaskan \omterm{def:g10:energy-conservation:energy}{energi} tepat ketika hasilnya duduk \emph{lebih tinggi} pada kurvanya; karena itu dua
siasat yang berlawanan sama-sama menguntungkan, dan keduanya berjalan menuju besi ---
\emph{membelah} inti yang terberat, dan \emph{menggabungkan} yang teringan.
\end{proposition}
\admitted

\begin{remark}[Mengapa ada puncak]\label{rem:g12:nuclear-energy:whypeak}
Tarik-menarik pada \cref{ch:g11:fundamental-interactions}: lem kuatnya mengikat
hanya tetangga yang bersentuhan, tolakan Coulombnya membentang seluruh intinya --- inti
yang kecil kebanyakan permukaan, yang besar kebanyakan tolakan, dan besi jalan tengahnya.
Pembukuan yang jujur bersemayam pada jilid universitasnya.
\end{remark}

\begin{omfigure}
\begin{tikzpicture}
\begin{axis}[omaxis, width=9.2cm, height=5.6cm,
    xlabel={nomor massa $A$}, ylabel={$E_b/A$ (\unit{MeV})},
    xmin=0, xmax=248, ymin=0, ymax=10.4, ytick={0,2,4,6,8}, xtick={4,56,120,180,235}]
  \addplot[omDef, thick, smooth] coordinates {(1,0) (2,1.1) (3,2.8) (6,5.3) (9,6.5)
    (12,7.7) (20,8.0) (40,8.6) (56,8.8) (90,8.7) (120,8.5) (150,8.3) (190,7.9) (235,7.6)};
  \addplot[only marks, mark=*, mark size=1.6pt, omThm] coordinates {(4,7.1) (56,8.8) (235,7.6)};
  \node[font=\small, right] at (axis cs:8,7.1) {$^{4}$He};
  \node[font=\small, above] at (axis cs:56,8.9) {$^{56}$Fe};
  \node[font=\small, above] at (axis cs:235,7.75) {$^{235}$U};
  \draw[-Stealth, omProp, thick] (axis cs:16,3.4) -- (axis cs:44,7.6)
    node[midway, below right, font=\small] {fusi};
  \draw[-Stealth, omProp, thick] (axis cs:222,6.9) -- (axis cs:130,7.8)
    node[midway, below, font=\small] {fisi};
\end{axis}
\end{tikzpicture}

{\small \omterm{def:g12:nuclear-energy:pernucleon}{Energi ikat tiap nukleon}: tanjakan yang curam, \omterm{prop:g12:nuclear-energy:ironpeak}{puncak besi}, luncuran yang panjang ---
\omterm{def:g12:nuclear-energy:fusion}{fusi} menanjak dari kiri, \omterm{def:g12:nuclear-energy:fission}{fisi} dari kanan, keduanya melepaskan \omterm{def:g10:energy-conservation:energy}{energi}.}
\end{omfigure}

\section{Fisi dan reaksi berantai}

\begin{definition}[Fisi]\label{def:g12:nuclear-energy:fission}
\emph{Fisi}\index{fisi} adalah pembelahan inti \omterm{def:g10:universal-gravitation:weight}{berat} menjadi dua
serpihan berukuran sedang. Uranium-235 bersifat \emph{fisil}\index{fisil}: sebutir
neutron lambat, yang terserap, meninggalkan intinya begitu terguncang sehingga ia terkoyak dua,
memuntahkan dua atau tiga neutron yang segar ---
${}^{1}_{0}\mathrm{n} + {}^{235}_{92}\mathrm{U} \to \text{dua serpihan} +
2\text{--}3\;{}^{1}_{0}\mathrm{n}$, dengan pasangannya beragam dari fisi ke fisi.
\end{definition}

\begin{example}[Satu fisi, yang ditimbang]\label{ex:g12:nuclear-energy:fissionbalance}
Salah satu jalur yang sering, menurut \cref{met:g12:nuclear-energy:balance}:
\[
{}^{1}_{0}\mathrm{n} + {}^{235}_{92}\mathrm{U} \longrightarrow
{}^{94}_{38}\mathrm{Sr} + {}^{140}_{54}\mathrm{Xe} + 2\,{}^{1}_{0}\mathrm{n}.
\]
Sebelumnya: $1.00866 + 234.99350 = \qty{236.00216}{u}$. Sesudahnya:
$93.89450 + 139.89200 + 2 \times 1.00866 = \qty{235.80382}{u}$. Karena itu
$\Delta m = \qty{0.19834}{u}$ dan $E \approx \qty{185}{MeV}$; dengan menghitung
peluruhan serpihan radioaktifnya kemudian (\cref{ch:g12:radioactive-decay}),
tiap \omterm{def:g12:nuclear-energy:fission}{fisinya} bernilai sekitar \qty{200}{MeV} --- beberapa \unit{eV} membeli satu \omterm{def:g11:fundamental-interactions:ladder}{atom}
batu bara.
\end{example}

\begin{definition}[Reaksi berantai dan massa kritis]\label{def:g12:nuclear-energy:chain}
Tiap \omterm{def:g12:nuclear-energy:fission}{fisi} dinyalakan satu neutron lalu membebaskan dua atau tiga: pada bongkahan
yang cukup besar mereka menghantam inti yang lain dan reaksinya menyuapi dirinya sendiri ---
sebuah \emph{reaksi berantai}\index{reaksi berantai}. Pada bongkahan yang kecil kebanyakan
neutronnya lolos lebih dahulu melalui permukaannya; kuantitas terkecil yang menopang rantainya
adalah \emph{massa kritis}\index{massa kritis} --- di bawahnya rantainya
padam, di atasnya tiap keturunannya berlipat: asas sebuah bom.
\end{definition}

\begin{omfigure}
\begin{tikzpicture}[font=\small, scale=0.9]
  \draw[-Stealth, thick] (-1.4,0) -- (-0.4,0) node[pos=0.15, above] {n};
  \fill[omDef!55] (0,0) circle (11pt); \node at (0,0) {$^{235}$U};
  \fill[omExo] (0.75,0.85) circle (4.5pt) node[above=3pt] {serpihan};
  \fill[omExo] (0.75,-0.85) circle (4.5pt);
  \foreach \y in {1.8,0,-1.8} {\draw[-Stealth] (0.5,0) -- (2.5,\y);
    \fill[omDef!55] (3,\y) circle (11pt); \node at (3,\y) {$^{235}$U};
    \foreach \a in {28,0,-28} \draw[-Stealth] (3.5,\y) -- ++(\a:1.1);}
\end{tikzpicture}

{\small Sebuah \omterm{def:g12:nuclear-energy:chain}{reaksi berantai}: satu neutron membelah satu inti, ketiga neutron
yang terbebas membelah tiga lagi --- sebuah reaktor hanya mengizinkan tepat satu untuk melanjutkan.}
\end{omfigure}

\begin{definition}[Reaktor yang dijinakkan]\label{def:g12:nuclear-energy:reactor}
Sebuah reaktor \omterm{def:g11:work-of-force:power}{daya} menahan rantainya tepat pada satu neutron tiap \omterm{def:g12:nuclear-energy:fission}{fisi}. Bahan bakarnya:
peleti uranium yang diperkaya sedikit dengan uranium-235.
\emph{Moderator}\index{moderator} (yang biasanya berupa air biasa) memperlambat
neutron \omterm{def:g12:nuclear-energy:fission}{fisi} yang cepat lewat tumbukan --- neutron yang lambat jauh lebih baik dalam memicu
\omterm{def:g12:nuclear-energy:fission}{fisi} berikutnya; \emph{batang kendali}\index{batang kendali} (boron, kadmium)
melahap neutronnya --- ketika didorong masuk, rantainya mati; ketika ditarik keluar, rantainya bergegas.
Guncangan serpihannya menjadi kalor, lalu uap, lalu listrik
(\cref{ch:g10:energy-conservation}).
\end{definition}

\begin{example}[Sebutir merica lawan sebuah truk]\label{ex:g12:nuclear-energy:gram}
Segram uranium-235 memegang
$N = \num{1.0e-3} / (235 \times \num{1.66054e-27}) \approx \num{2.56e21}$
inti; pada $\qty{200}{MeV} = \qty{3.2e-11}{J}$ masing-masing, membelah semuanya
menghasilkan sekitar \qty{8.2e10}{J} --- \omterm{def:g10:energy-conservation:forms}{energi kimia} $\num{8.2e10} /
\num{3.0e7} \approx \qty{2.7e3}{kg}$ batu bara: sebutir merica lawan tiga ton.
\end{example}

\section{Fusi: api milik Matahari sendiri}

\begin{definition}[Fusi]\label{def:g12:nuclear-energy:fusion}
\emph{Fusi}\index{fusi} adalah penggabungan dua inti ringan menjadi satu inti yang lebih berat
dan lebih erat terikat. Reaksi yang paling terjangkau menikahkan deuterium dan
tritium, \omterm{def:g11:nucleus-radioactivity:notation}{isotop} \omterm{def:g10:universal-gravitation:weight}{berat} hidrogen (\cref{ch:g11:nucleus-radioactivity}):
\[
{}^{2}_{1}\mathrm{H} + {}^{3}_{1}\mathrm{H} \longrightarrow
{}^{4}_{2}\mathrm{He} + {}^{1}_{0}\mathrm{n}.
\]
\end{definition}

\begin{example}[D--T, yang ditimbang]\label{ex:g12:nuclear-energy:dt}
Sebelumnya: $2.01355 + 3.01550 = \qty{5.02905}{u}$; sesudahnya:
$4.00151 + 1.00866 = \qty{5.01017}{u}$; jadi $\Delta m = \qty{0.01888}{u}$ dan
$E \approx \qty{17.6}{MeV}$, empat perlimanya pada neutronnya --- kurang daripada
sebuah \omterm{def:g12:nuclear-energy:fission}{fisi}, tetapi dari lima \omterm{def:g11:fundamental-interactions:ladder}{nukleon} dan bukan $236$: sekitar \qty{3.4e14}{J} tiap
\omterm{def:g10:orders-of-magnitude:unit}{kilogram} bahan bakarnya, empat kali \omterm{def:g12:nuclear-energy:fission}{fisinya}. Helium-4
(\cref{ex:g12:nuclear-energy:helium}) sekali lagi menjadi abunya.
\end{example}

\begin{definition}[Penghalang Coulomb]\label{def:g12:nuclear-energy:barrier}
Kedua intinya sama-sama positif: untuk bersentuhan, keduanya harus lebih dahulu menanjaki bukit
tolakan listriknya (\cref{ch:g11:electric-gravitational-fields}) --- yaitu
\emph{penghalang Coulomb}\index{penghalang Coulomb}. Hanya di dekat \qty{1e8}{K}
tumbukannya membawa \omterm{def:g10:energy-conservation:energy}{energi} untuk melintasinya; bahan bakarnya lalu menjadi
\emph{plasma}\index{plasma} berisi inti telanjang dan elektron. \omterm{def:g12:nuclear-energy:fission}{Fisi} tidak membayar
pungutan semacam itu: pemicunya, yaitu neutron, bermuatan netral, dan bekerja dalam keadaan dingin.
\end{definition}

\begin{example}[Matahari berjalan dengannya]\label{ex:g12:nuclear-energy:sun}
Teras Matahari, pada $\qty{1.5e7}{K}$, memadukan hidrogen biasa selangkah demi selangkah ---
yaitu \emph{rantai proton--proton}, yang berdampak bersih
$4\,{}^{1}_{1}\mathrm{H} \to {}^{4}_{2}\mathrm{He} + 2\,{}^{0}_{+1}\mathrm{e}
+ \text{radiasi}$, sekitar \qty{26}{MeV} tiap helium: $0.7\%$ massa
hidrogennya pergi sebagai \omterm{def:g10:energy-conservation:energy}{energi}. Cahaya matahari yang \omterm{def:g12:sound-acoustics:timbre}{spektrumnya} kita baca pada
\cref{ch:g10:light-spectra} adalah massa ini, yang tiba terlambat delapan menit; setiap
bintang bersinar dengan menyusuri kurvanya menuju besi.
\end{example}

\begin{remark}[ITER: membotolkan sebuah bintang]\label{rem:g12:nuclear-energy:iter}
Tidak ada padatan yang dapat memegang \omterm{def:g12:nuclear-energy:barrier}{plasma} pada \qty{1.5e8}{K}; sebuah \emph{tokamak} mengurungnya di dalam
\omterm{def:g11:magnetic-fields:bfield}{medan magnet} (\cref{ch:g11:magnetic-fields}), sebuah donat \omterm{def:g11:electric-gravitational-fields:lines}{garis medan} yang
zarah bermuatannya menyusuri secara membelit tanpa menyentuh dindingnya; ITER, yang dibangun untuk
melepaskan sepuluh kali \omterm{def:g11:work-of-force:power}{daya} pemanasnya, adalah langkah yang sekarang. Patahkanlah
pengurungannya dan \omterm{def:g12:nuclear-energy:barrier}{plasmanya} mendingin lalu berhenti dalam beberapa sekon: tanpa \omterm{def:g12:nuclear-energy:chain}{reaksi berantai}, tanpa
\omterm{def:g12:nuclear-energy:chain}{massa kritis} --- yang sulit bukanlah memadamkan apinya melainkan menjaganya tetap menyala.
\end{remark}

\begin{example}[Tangga energinya, tiap kilogram]\label{ex:g12:nuclear-energy:ladder}
\omterm{def:g10:energy-conservation:energy}{Energi} tiap \omterm{def:g10:orders-of-magnitude:unit}{kilogram} bahan bakar: batu bara, \qty{3.0e7}{J}; uranium-235 lewat \omterm{def:g12:nuclear-energy:fission}{fisi},
\qty{8.2e13}{J}; deuterium--tritium lewat \omterm{def:g12:nuclear-energy:fusion}{fusi}, \qty{3.4e14}{J}; pengubahan
seluruhnya ($E = mc^2$), \qty{9.0e16}{J}. Enam orde besaran memisahkan
kimia dari intinya --- itulah sebabnya sebuah reaktor diisi ulang dengan truk
dan pembangkit batu bara dengan rangkaian kereta.
\end{example}

\section{Latihan}

\begin{exercise}[$\star$]\label{exo:g12:nuclear-energy:1}
Ubahlah: (a) \qty{1}{eV} dan \qty{1}{MeV} menjadi \omterm{def:g11:work-of-force:work}{joule}; (b) \qty{200}{MeV} milik
satu \omterm{def:g12:nuclear-energy:fission}{fisi} menjadi \omterm{def:g11:work-of-force:work}{joule}; (c) periksalah dari $\qty{1}{u} = \qty{1.66054e-27}{kg}$
bahwa setara \omterm{def:g10:energy-conservation:energy}{energinya} dekat dengan \qty{931.5}{MeV}.
\end{exercise}

\begin{exercise}[$\star$]\label{exo:g12:nuclear-energy:2}
Sebuah dadu gula bermassa \qty{6.0}{g}: hitunglah nilai $E = mc^2$ penuhnya. Sebuah
rumah tangga menarik $\qty{10}{kW.h} = \qty{3.6e7}{J}$ tiap hari --- selama berapa
tahun dadu itu, jika diubah seluruhnya, dapat menyalakannya?
\end{exercise}

\begin{exercise}[$\star$]\label{exo:g12:nuclear-energy:3}
Inti karbon-12 bermassa \qty{11.99671}{u}: hitunglah \omterm{def:g12:nuclear-energy:defect}{defek massanya},
\omterm{def:g12:nuclear-energy:defect}{energi ikatnya} dalam \unit{MeV}, dan \omterm{def:g12:nuclear-energy:pernucleon}{energi ikat tiap nukleonnya}
($m_p = \qty{1.00728}{u}$, $m_n = \qty{1.00866}{u}$).
\end{exercise}

\begin{exercise}[$\star$]\label{exo:g12:nuclear-energy:4}
Lengkapilah tiap \omterm{def:g12:nuclear-energy:fission}{fisinya}, dengan menyebutkan $A$, $Z$ dan unsurnya ($Z = 36$: kripton,
$Z = 52$: telurium):
(a) ${}^{1}_{0}\mathrm{n} + {}^{235}_{92}\mathrm{U} \to
{}^{144}_{56}\mathrm{Ba} + {}^{A}_{Z}\mathrm{X} + 3\,{}^{1}_{0}\mathrm{n}$;
(b) ${}^{1}_{0}\mathrm{n} + {}^{235}_{92}\mathrm{U} \to
{}^{97}_{40}\mathrm{Zr} + {}^{A}_{Z}\mathrm{Y} + 2\,{}^{1}_{0}\mathrm{n}$.
\end{exercise}

\begin{exercise}[$\star$]\label{exo:g12:nuclear-energy:5}
Bacalah kurva pada \cref{prop:g12:nuclear-energy:ironpeak}: taksirlah $E_b/A$ untuk
helium-4, karbon-12, besi-56 dan uranium-235. Mana yang paling erat terikat?
Mana yang dapat melepaskan \omterm{def:g10:energy-conservation:energy}{energi} lewat \omterm{def:g12:nuclear-energy:fission}{fisi}, dan mana lewat \omterm{def:g12:nuclear-energy:fusion}{fusi}?
\end{exercise}

\begin{exercise}[$\star\star$]\label{exo:g12:nuclear-energy:6}
Deuteron ${}^{2}_{1}\mathrm{H}$ bermassa \qty{2.01355}{u}: hitunglah
\omterm{def:g12:nuclear-energy:defect}{energi ikatnya} dalam \unit{MeV}. Ikatan kimia memegang beberapa \unit{eV} --- dengan
faktor berapa ``ikatan'' intinya mengalahkannya?
\end{exercise}

\begin{exercise}[$\star\star$]\label{exo:g12:nuclear-energy:7}
Timbanglah \omterm{def:g12:nuclear-energy:fission}{fisi} ${}^{1}_{0}\mathrm{n} + {}^{235}_{92}\mathrm{U} \to
{}^{144}_{56}\mathrm{Ba} + {}^{89}_{36}\mathrm{Kr} + 3\,{}^{1}_{0}\mathrm{n}$:
massanya \qty{234.99350}{u} (U), \qty{143.89220}{u} (Ba), \qty{88.89790}{u}
(Kr), \qty{1.00866}{u} (n). Berapa \omterm{def:g10:energy-conservation:energy}{energi} yang dilepaskan, dalam \unit{MeV} dan dalam \omterm{def:g11:work-of-force:work}{joule}?
\end{exercise}

\begin{exercise}[$\star\star$]\label{exo:g12:nuclear-energy:8}
\omterm{def:g12:nuclear-energy:fusion}{Fusi} yang lain: ${}^{2}_{1}\mathrm{H} + {}^{2}_{1}\mathrm{H} \to
{}^{3}_{2}\mathrm{He} + {}^{1}_{0}\mathrm{n}$, dengan
$m({}^{3}_{2}\mathrm{He}) = \qty{3.01493}{u}$. Hitunglah \omterm{def:g10:energy-conservation:energy}{energi} yang dilepaskannya.
Mengapa jauh lebih kecil daripada D--T? (Timbanglah letak tiap hasilnya pada kurvanya.)
\end{exercise}

\begin{exercise}[$\star\star$]\label{exo:g12:nuclear-energy:9}
Jelaskan, beberapa baris masing-masing: (a) mengapa \omterm{def:g12:nuclear-energy:fission}{fisi} tidak memerlukan pemanasan sementara \omterm{def:g12:nuclear-energy:fusion}{fusi}
menuntut sekitar \qty{1e8}{K}; (b) bagaimana Matahari sanggup pada ``hanya'' \qty{1.5e7}{K}
(ukuran kerumunannya, anggotanya yang tercepat); (c) mengapa pembangkit \omterm{def:g12:nuclear-energy:fusion}{fusi} tidak dapat
lepas kendali seperti \omterm{def:g12:nuclear-energy:chain}{reaksi berantai}.
\end{exercise}

\begin{exercise}[$\star\star$]\label{exo:g12:nuclear-energy:10}
Reaksi bersih Matahari mengubah empat proton ($m_p = \qty{1.00728}{u}$) menjadi satu
inti helium-4 (\qty{4.00151}{u}): hitunglah \omterm{def:g10:energy-conservation:energy}{energi} tiap helium yang terbuat, dalam
\unit{MeV}, dan bagian massa awalnya yang terubah.
\end{exercise}

\begin{exercise}[$\star\star$]\label{exo:g12:nuclear-energy:11}
Pada reaktor bermoderator air: (a) apa yang dikerjakan \omterm{def:g12:nuclear-energy:reactor}{moderatornya}, dan mengapa
rantainya membutuhkannya? (b) \omterm{def:g12:nuclear-energy:reactor}{batang kendalinya}? (c) air pendinginnya \emph{adalah}
\omterm{def:g12:nuclear-energy:reactor}{moderatornya}: mengapa kehilangannya cenderung mencekik rantainya dan bukan mempercepatnya?
\end{exercise}

\begin{exercise}[$\star\star\star$]\label{exo:g12:nuclear-energy:12}
Sebuah rantai yang tak terkendali berlipat dua tiap keturunan. Dari satu \omterm{def:g12:nuclear-energy:fission}{fisi}, kira-kira berapa
keturunan sampai $\num{2.56e21}$ inti segram uranium-235
terbelah ($2^{10} \approx 10^{3}$)? Pada sekitar \qty{10}{ns} tiap keturunan,
berapa lama itu --- dan mengapa, karenanya, kekritisannya dijaga begitu saksama?
\end{exercise}

\begin{exercise}[$\star\star\star$]\label{exo:g12:nuclear-energy:13}
Dapatkah kita menambang \omterm{def:g10:energy-conservation:energy}{energi} dengan membelah besi? Timbanglah
${}^{56}_{26}\mathrm{Fe} \to 2\,{}^{28}_{14}\mathrm{Si}$, dengan
$m(\mathrm{Fe}) = \qty{55.92066}{u}$, $m(\mathrm{Si}) = \qty{27.96924}{u}$.
Hitunglah $\Delta m$ lalu simpulkanlah, dari kurvanya, mengapa besi adalah
abu inti: bahan bakar bagi apa pun tidak.
\end{exercise}

\begin{exercise}[$\star\star\star$]\label{exo:g12:nuclear-energy:14}
Dari $m({}^{56}_{26}\mathrm{Fe}) = \qty{55.92066}{u}$ dan
$m({}^{235}_{92}\mathrm{U}) = \qty{234.99350}{u}$, hitunglah $E_b/A$ untuk keduanya
($m_p = \qty{1.00728}{u}$, $m_n = \qty{1.00866}{u}$). Serpihan \omterm{def:g12:nuclear-energy:fission}{fisinya} duduk
di dekat \qty{8.5}{MeV} tiap \omterm{def:g11:fundamental-interactions:ladder}{nukleon}: \emph{taksirlah} dari kedua arasnya \omterm{def:g10:energy-conservation:energy}{energi}
satu \omterm{def:g12:nuclear-energy:fission}{fisi}, lalu bandingkan dengan \cref{ex:g12:nuclear-energy:fissionbalance}.
\end{exercise}

\begin{exercise}[$\star\star\star$]\label{exo:g12:nuclear-energy:15}
ITER membidik \qty{500}{MW} \omterm{def:g11:work-of-force:power}{daya} \omterm{def:g12:nuclear-energy:fusion}{fusi}. Pada \qty{17.6}{MeV} tiap reaksi D--T
(\qty{5.029}{u} bahan bakar yang terpakai), carilah banyaknya reaksi tiap sekon,
lalu bahan bakar yang terbakar tiap hari. Sebuah tungku batu bara \qty{500}{MW} membakar sekitar
\qty{1.4e6}{kg} tiap hari: bandingkanlah.
\end{exercise}

\section{Soal: Menyalakan sebuah kota}

\begin{problem}[{Soal akhir pekan --- menyalakan sebuah kota: satu gigawatt yang mantap bagi
setengah juta orang, yang dibeli dengan empat cara --- peleti uranium, rangkaian kereta
batu bara, sekolam air laut, zat milik Matahari sendiri --- dengan satu kilogram
massa yang lenyap yang sama bersembunyi di bawah keempatnya}]
\label{pb:g12:nuclear-energy:1}
Sebuah kota berpenduduk $500\,000$ jiwa menarik \omterm{prop:g11:circuits-and-power:power}{daya listrik} yang mantap
$P = \qty{1.0}{GW}$; pembangkitnya, apa pun bahan bakarnya, mengubah kalor menjadi
listrik dengan dayaguna $33\%$ (\cref{ch:g10:energy-conservation}).
Datanya: satu \omterm{def:g12:nuclear-energy:fission}{fisi} uranium-235, \qty{200}{MeV} seluruhnya; batu bara, \qty{3.0e7}{J/kg};
satu \omterm{def:g12:nuclear-energy:fusion}{fusi} D--T, \qty{17.6}{MeV} dari \qty{5.029}{u} bahan bakar;
$\qty{1}{u} = \qty{1.66054e-27}{kg}$; $\qty{1}{eV} = \qty{1.60e-19}{J}$; satu
tahun $= \qty{3.16e7}{s}$; Matahari: \qty{3.9e26}{W}, \qty{2.0e30}{kg}; semeter
kubik air laut memegang sekitar \qty{33}{g} deuterium.

\medskip
\noindent\textbf{Bagian I --- Pembangkit \omterm{def:g12:nuclear-energy:fission}{fisinya}.}
\begin{enumerate}
  \item \omterm{def:g10:energy-conservation:power}{Daya} kalor berapa yang harus diserahkan reaktornya?
  \item Nyatakan \omterm{def:g10:energy-conservation:energy}{energi} satu \omterm{def:g12:nuclear-energy:fission}{fisi} dalam \omterm{def:g11:work-of-force:work}{joule}.
  \item Berapa \omterm{def:g12:nuclear-energy:fission}{fisi} tiap sekon yang menjalankan kotanya?
  \item Berapa massa satu inti uranium-235 ($A = 235$), lalu uranium-235
        yang terpakai tiap sekon?
  \item Perolehlah uranium-235 yang terpakai dalam setahun, dalam \omterm{def:g10:orders-of-magnitude:unit}{kilogram}.
  \item Bahan bakarnya diperkaya sampai $4.0\%$ uranium-235: massa berapa yang
        dituntut setahun, dan muatkah pada satu truk?
\end{enumerate}

\medskip
\noindent\textbf{Bagian II --- Buku besar batu baranya.}
\begin{enumerate}[resume]
  \item Hitunglah \omterm{def:g10:energy-conservation:energy}{energi} kalor yang dikonsumsi kotanya dalam setahun.
  \item Massa batu bara berapa yang menyerahkannya?
  \item Hitunglah nisbah (massa batu bara)/(massa uranium-235) untuk setahun, lalu
        periksalah terhadap tangga pada \cref{ex:g12:nuclear-energy:ladder}.
  \item Sebuah kereta barang menarik \qty{3.0e6}{kg}: berapa kereta tiap tahun, dan
        tiap hari? Satu kalimat: bayangkanlah kedua jalur pasokannya.
\end{enumerate}

\medskip
\noindent\textbf{Bagian III --- Impian \omterm{def:g12:nuclear-energy:fusion}{fusinya}.}
\begin{enumerate}[resume]
  \item Nyatakan \omterm{def:g10:energy-conservation:energy}{energi} satu \omterm{def:g12:nuclear-energy:fusion}{fusi} D--T dalam \omterm{def:g11:work-of-force:work}{joule}.
  \item Hitunglah \omterm{def:g10:energy-conservation:energy}{energi} yang dilepaskan tiap \omterm{def:g10:orders-of-magnitude:unit}{kilogram} bahan bakar D--T; bandingkan dengan
        tangganya.
  \item Massa bahan bakar D--T berapa yang menutupi setahun kotanya? Berapa massa darinya yang berupa
        deuterium ($2.014$ dari \qty{5.029}{u} itu)?
  \item Volume air laut berapa yang memuat deuterium itu? Bandingkan dengan
        kolam Olimpiade, sekitar \qty{2.5e3}{m^3}. (Tritiumnya dibiakkan dari litium.)
  \item Mengapa pembangkit ini tidak membawa \omterm{def:g12:nuclear-energy:chain}{massa kritis} dan \omterm{def:g12:nuclear-energy:chain}{reaksi berantai} yang
        harus ditakuti? Lalu, apa bagian yang sulitnya?
\end{enumerate}

\medskip
\noindent\textbf{Bagian IV --- Bintang yang membayar dengan massa.}
\begin{enumerate}[resume]
  \item Dari $E = mc^2$, massa berapa yang diubah Matahari tiap sekon?
  \item Massa berapa tiap sekon, dan tiap tahun, yang harus diubah sumber
        \qty{3.0}{GW} kalor \emph{mana pun}? Periksalah bahwa ketiga bahan bakarnya
        menyerahkan massa yang sama ini.
  \item Massa seluruhnya berapa yang sudah diradiasikan Matahari selama $\num{4.5e9}$
        tahunnya, dan berapa bagian Matahari nilai itu?
  \item \omterm{def:g12:nuclear-energy:fusion}{Fusi} dapat menyadap sekitar $0.07\%$ massa Matahari (bagian terasnya yang dapat
        terbakar): taksirlah umur seluruhnya sebagai bintang, dan apa yang tersisa.
  \item Tutuplah pemeriksaannya dalam tiga kalimat: setahun kotanya dalam ton
        batu bara, \omterm{def:g10:orders-of-magnitude:unit}{kilogram} uranium-235 dan \omterm{def:g10:orders-of-magnitude:unit}{kilogram} bahan bakar D--T --- dan
        satu \omterm{def:g10:orders-of-magnitude:unit}{kilogram} massa yang, pada setiap kisahnya, menjadi cahaya kotanya.
\end{enumerate}
\end{problem}
